% ===========================================================
%  《线性代数9讲》第 1 讲 行列式 —— 片段 B
%  覆盖 PDF p17–p20（印刷页 6–9）；四页均已逐页转录
%  开头顶格承接 PDF p16 例 1.4【注】(2)（数学归纳法）的续页
%  约定：原稿蓝色 OPD 方法标注（D_1 / D_21 / D_22 / D_23）一律以
%        "% OPD" 注释保留，未排入正文（依 AGENTS.md §7.1 第 2 条）；
%        其余蓝色旁注（交叉引用、用好公式）照原文排入。
%  ⚠️ 主控已改：\fsec 标题中的 OPD 标记（O_2（盯住目标…））已由主控剥离，
%     只保留数学分类（具体型／抽象型）。请勿"按原文恢复"该标记。
%     本片转录已于 21:07 结束，文件此后由主控接管。
% ===========================================================

% -----------------------------------------------------------
% PDF p17（印刷页 6）
% -----------------------------------------------------------
% 原稿此处为青色框，系 p16 例 1.4【注】(2) 的续页
\begin{fnote}
% OPD 蓝色标注（原稿）：设 $D_{21}$（观察研究对象）→ 指向下式 $D_k$
设 $D_k=b^{k}+a_1b^{k-1}+\cdots+a_{k-1}b+a_k$, ①

现在来看 $D_{k+1}$，将 $D_{k+1}$ 按第 1 列展开，由上述解析，知
\fml{D_{k+1}=bD_k+a_{k+1},\qquad\text{②}}
将归纳假设①式代入②式，得
\fml{D_{k+1}=b\,(b^{k}+a_1b^{k-1}+\cdots+a_{k-1}b+a_k)+a_{k+1}}
\fml{=b^{k+1}+a_1b^{k}+\cdots+a_{k-1}b^{2}+a_kb+a_{k+1},}
因此①式对任何正整数 $k$ 都成立，即得
\fml{D_k=b^{k}+a_1b^{k-1}+\cdots+a_{k-1}b+a_k.}
\end{fnote}

\begin{fcard}{例 1.5}
证明：$n$ 阶行列式
\fml{D_n=\begin{vmatrix}
2a & 1 & 0 & \cdots & 0 & 0\\
a^{2} & 2a & 1 & \cdots & 0 & 0\\
0 & a^{2} & 2a & \cdots & 0 & 0\\
\vdots & \vdots & \vdots & & \vdots & \vdots\\
0 & 0 & 0 & \cdots & 2a & 1\\
0 & 0 & 0 & \cdots & a^{2} & 2a
\end{vmatrix}=(n+1)a^{n}.}

【证】用第二数学归纳法.

当 $n=1$ 时，$D_1=2a=(1+1)a^{1}$，命题成立.

当 $n=2$ 时，$D_2=\begin{vmatrix}2a & 1\\ a^{2} & 2a\end{vmatrix}=4a^{2}-a^{2}=3a^{2}=(2+1)a^{2}$，命题成立.

假设当 $n<k$ 时，命题成立，则当 $n=k$（$\geqslant 3$）时，$D_k$ 按第 1 列展开，得
\fml{D_k=2aD_{k-1}+(-1)^{1+2}a^{2}\begin{vmatrix}
1 & 0 & 0 & \cdots & 0 & 0\\
a^{2} & 2a & 1 & \cdots & 0 & 0\\
0 & a^{2} & 2a & \cdots & 0 & 0\\
\vdots & \vdots & \vdots & & \vdots & \vdots\\
0 & 0 & 0 & \cdots & 2a & 1\\
0 & 0 & 0 & \cdots & a^{2} & 2a
\end{vmatrix}_{k-1}}
\fml{=2aD_{k-1}-a^{2}D_{k-2}}
\fml{=2a\,(k+1-1)\,a^{k-1}-a^{2}(k-2+1)\,a^{k-2}=(k+1)a^{k},}
得证，命题成立.
\end{fcard}

\begin{fnote}
【注】一般来说，当命题直接要求计算行列式时，优先考虑递推法，如例 1.4；当命题给出行列式的结果，要求证明时，优先考虑数学归纳法，如例 1.5.
\end{fnote}

% -----------------------------------------------------------
% PDF p18（印刷页 7）
% -----------------------------------------------------------
\fsec{二、抽象型行列式的计算：$a_{ij}$ 未给出}

\begin{fcard}{1. 用行列式的性质}
用行列式的性质将所求行列式进一步化成已知行列式.
\end{fcard}

\begin{fcard}{2. 用矩阵知识}
（1）设 $C=AB$，$A$，$B$ 为同阶方阵，则 $|C|=|AB|=|A||B|$.
\end{fcard}

% OPD 蓝色标注（原稿）：→$D_{23}$（化归经典形式）
\begin{fnote}
【注】要善于发现 $AB$，如 $A=[\alpha_1,\alpha_2,\alpha_3]$，$B=\begin{bmatrix}a & b & b\\ b & a & b\\ b & b & a\end{bmatrix}$，$C=\begin{bmatrix}1 & 2 & 3\\ 2 & 3 & 4\\ 3 & 4 & 5\end{bmatrix}$.

令
\fml{D_1=[a\alpha_1+b(\alpha_2+\alpha_3),\,a\alpha_2+b(\alpha_1+\alpha_3),\,a\alpha_3+b(\alpha_1+\alpha_2)]=AB;}
\fml{D_2=[\alpha_1+2\alpha_2+3\alpha_3,\,2\alpha_1+3\alpha_2+4\alpha_3,\,3\alpha_1+4\alpha_2+5\alpha_3]=AC.}
这样利于使用 $|D_1|=|A||B|$，$|D_2|=|A||C|$.
\end{fnote}

（2）设 $C=A+B$，$A$，$B$ 为同阶方阵，则 $|C|=|A+B|$，但由于 $|A+B|$ 不一定等于 $|A|+|B|$，故需对 $|A+B|$ 作恒等变形，转化为矩阵乘积的行列式. 这里的恒等变形一般是：①由题设条件，如 $E=AA^{T}$；②用 $E=AA^{-1}$ 等.
% OPD 蓝色标注（原稿）：→$D_{23}$（化归经典形式）  （指向「转化为矩阵乘积的行列式」）
% OPD 蓝色标注（原稿）：→$D_{22}$（转换等价表述）  （指向「①由题设条件，如 $E=AA^{T}$」）

\begin{fcard}{例 1.6}
设 $A$，$B$ 为 $n$ 阶正交矩阵，则 $n$ 为奇数时，$|(A-B)(A+B)|=\underline{\qquad}$.

【分析】$AA^{T}=E(A^{T}=A^{-1})$；$BB^{T}=E(B^{T}=B^{-1})$.

【解】应填 0.
% OPD 蓝色标注（原稿）：=$E$（$D_{22}$（转换等价表述））
\fml{|(A-B)(A+B)|=|(A-B)^{T}(A+B)|=|A^{T}A-B^{T}A+A^{T}B-B^{T}B|=|A^{T}B-B^{T}A|,}
又 $|(A-B)(A+B)|=|(A+B)^{T}(A-B)|=(-1)^{n}|A^{T}B-B^{T}A|$，故当 $n$ 为奇数时，$|(A-B)(A+B)|=0$.
\end{fcard}

\begin{fnote}
【注】$|A|=0$ 可由以下条件推知.

①$|A|=k|A|$，$k\neq 1$，常考 $|A|=(-1)^{n}|A|$($n$ 为奇数).

②$A$ 有特征值 0.

③$A_{n\times n}x=\mathbf{0}$ 有非零解.

④$r(A_{m\times n})<n$.

⑤$A$ 的极大无关组成员个数 $<n$（也即 $A$ 的行（列）向量组线性相关）.

⑥实在不行，或显而易见时，立即推——放弃考研！不，用反证法，设 $A$ 可逆，即 $|A|\neq 0$，推得矛盾，即可得证！
\end{fnote}

% -----------------------------------------------------------
% PDF p19（印刷页 8）
% -----------------------------------------------------------
% OPD 蓝色标注（原稿）：→$D_1$（常规操作）$+D_{23}$（化归经典形式）
（3）设 $A$ 为 $n$ 阶矩阵，则 $|A^{*}|=|A|^{n-1}$，$|(A^{*})^{*}|=||A|^{n-2}A|=|A|^{(n-1)^{2}}$. 更全面的公式总结在第 3 讲的“四、2.（2）”处.

\begin{fcard}{例 1.7}
设 $A$ 是 4 阶矩阵，且其伴随矩阵 $A^{*}$ 的特征值为 1，$-1$，$-2$，4，则 $|A^{3}+2A^{2}-A-3E|=\underline{\qquad}$.

【解】应填 $-\dfrac{253}{8}$.

% 原稿此处右侧另有一张 3 行 3 列表格
\fml{\begin{array}{|c|c|c|}\hline
A & A^{*} & f(A)\\ \hline
\lambda & \dfrac{|A|}{\lambda} & f(\lambda)\\ \hline
\xi & \xi & \xi\\ \hline
\end{array}}

由于 $|A^{*}|=1\times(-1)\times(-2)\times4=8\neq 0$，可知 $A^{*}$ 可逆，于是 $A$ 可逆. 又 $|A^{*}|=|A|^{n-1}=|A|^{3}=8$，得 $|A|=2$. 故 $A$ 的特征值 $\lambda_A=\dfrac{|A|}{\lambda_{A^{*}}}$，即为 $2$，$-2$，$-1$，$\dfrac{1}{2}$.

设 $f(A)=A^{3}+2A^{2}-A-3E$，则 $f(A)$ 的特征值为

% 原稿蓝色旁注（交叉引用）：见第 7 讲的“一、（3）①”中的表格：$f(A)$ 的特征值为$f(\lambda)$
见第 7 讲的“一、（3）①”中的表格：$f(A)$ 的特征值为 $f(\lambda)$.

\fml{f(2)=2^{3}+2\times2^{2}-2-3=11,}
\fml{f(-2)=(-2)^{3}+2\times(-2)^{2}+2-3=-1,}
\fml{f(-1)=-1+2+1-3=-1,}
\fml{f\!\left(\frac{1}{2}\right)=\left(\frac{1}{2}\right)^{3}+2\times\left(\frac{1}{2}\right)^{2}-\frac{1}{2}-3=-\frac{23}{8},}
故
\fml{|A^{3}+2A^{2}-A-3E|=f(2)\cdot f(-2)\cdot f(-1)\cdot f\!\left(\frac{1}{2}\right)=-\frac{253}{8}.}
\end{fcard}

\begin{fcard}{3. 用相似理论}
% OPD 蓝色标注（原稿）：→$D_1$（常规操作）$+D_{22}$（转换等价表述）
若 $A$ 相似于 $B$，则 $|A|=|B|$，且 $|A|=\prod\limits_{i=1}^{n}\lambda_i$，主要考：

（1）用行列式 $|\lambda E-A|=0$ 求特征值；

（2）用特征值求行列式.
\end{fcard}

\begin{fcard}{例 1.8}
设 $A=k\begin{bmatrix}
1 & a & a & \cdots & a\\
a & 1 & a & \cdots & a\\
\vdots & \vdots & \vdots & & \vdots\\
a & a & a & \cdots & 1
\end{bmatrix}_{n\times n}$，$0<a\leqslant 1$，$k>0$，求 $A$ 的最大特征值.

% ---------------------------------------------------------
% PDF p20（印刷页 9）
% ---------------------------------------------------------
【解】
\fmll{|\lambda E-A|=\begin{vmatrix}
\lambda-k & -ka & \cdots & -ka\\
-ka & \lambda-k & \cdots & -ka\\
\vdots & \vdots & & \vdots\\
-ka & -ka & \cdots & \lambda-k
\end{vmatrix}=[\lambda-k+(n-1)(-ka)](\lambda-k+ka)^{n-1}=0,}

% 原稿蓝色旁注（带箭头指向上式结果）：用好公式：
用好公式：$\begin{bmatrix}a & b & \cdots & b\\ b & a & \cdots & b\\ \vdots & \vdots & & \vdots\\ b & b & \cdots & a\end{bmatrix}=[a+(n-1)b](a-b)^{n-1}$

故 $\lambda_1=k+(n-1)ka=k[1+(n-1)a]$（单重），$\lambda_2=k-ka=k(1-a)$（$n-1$ 重）.

由 $1+(n-1)a>1-a$，$k>0$，故最大特征值为 $k[1+(n-1)a]$.
\end{fcard}
